\documentclass[11pt]{article}
\usepackage[a4paper,left=2.5cm,right=2.5cm,top=2.2cm,bottom=2.2cm]{geometry}
\usepackage{amsmath,amssymb,amsthm,mathtools}
\usepackage{booktabs,enumitem}
\setlength{\emergencystretch}{2em}
\usepackage{xcolor}
\usepackage[colorlinks=true,linkcolor=blue!50!black,citecolor=blue!50!black,urlcolor=blue!50!black]{hyperref}
\usepackage[capitalize]{cleveref}

\newtheorem{theorem}{Theorem}
\newtheorem{proposition}[theorem]{Proposition}
\newtheorem{lemma}[theorem]{Lemma}
\newtheorem{corollary}[theorem]{Corollary}
\theoremstyle{definition}
\newtheorem{definition}[theorem]{Definition}
\newtheorem{remark}[theorem]{Remark}

\newcommand{\Tq}{T_{\star}}
\newcommand{\tr}{\operatorname{Tr}}
\newcommand{\re}{\operatorname{Re}}
\newcommand{\qone}{\mathcal{Q}_1}
\newcommand{\qab}{\mathcal{Q}_{1+AB}}
\newcommand{\PSD}{\succeq}
\hypersetup{pdftitle={The optimal three-outcome CGLMP value in the Gupta--Pienaar qplex model},
 pdfauthor={Yusuf Kaan Öğüt, Claude Opus 5.5 (AI), OpenAI Codex (GPT-6.2 Astra, AI)}}
\newcommand{\NBOX}{858}
\title{The optimal three-outcome CGLMP value\\in the Gupta--Pienaar qplex model}
\author{Yusuf Kaan \"{O}\u{g}\"{u}t\\[3pt]
{\small Claude Opus 5.5 (Anthropic) and OpenAI Codex (GPT-6.2 Astra)}\\[-1pt]
{\footnotesize AI contributors}}
\date{Revised 25 November 2026}
\begin{document}
\maketitle

\begin{abstract}
Gupta and Pienaar exhibited a three-outcome Bell correlation with CGLMP score
$2+2/\sqrt3$ in their $d=3$ qplex model and conjectured that it is optimal.
We prove the upper bound for all priors and weights. Equality forces Alice's outcome weights to be uniform for both settings.
For uniform weights and a flat prior, a Fourier reduction gives a four-dimensional
projector certificate. A Gram-matrix construction places every such correlation in
the macroscopically local set $\qone$ and gives the exact value
$\max_{\qone} I_3=2+2/\sqrt3$. The maximising example lies outside the almost-quantum set.
For biased weights we exhibit a rational qplex correlation and an integer matrix
certifying that it lies outside $\qone$. The general upper bound follows from a separate
computer-assisted proof. Fixed dual witnesses cover the four-dimensional weight domain,
including a neighbourhood of the optimum where the dual multipliers vary quadratically.
The accompanying verifier checks the witnesses in exact arithmetic using the Python standard library.
\end{abstract}

\section{Introduction}

QBism reads quantum states as an agent's probability assignments and the Born rule as a normative
coherence condition linking them~\cite{FMS2014,FuchsSchack2013}. Using a symmetric informationally complete (SIC)
reference measurement~\cite{Renes2004}, a quantum state of a $d$-level system becomes a probability vector on
$d^2$ outcomes, and the Born rule becomes the \emph{Urgleichung}~\cite{FuchsSchack2013}. The \emph{qplex}
programme~\cite{Qplex2017} asks which geometric conditions on such probability vectors single out quantum theory.
A qplex is a maximal set of probability vectors whose pairwise inner products lie in a fixed interval.
Quantum state space in SIC coordinates is one example, the Hilbert qplex.

Gupta and Pienaar~\cite{GuptaPienaar2026} studied bipartite correlations generated within this framework from the
perspective of one agent (``Bob''), who conditions on the other agent's (``Alice's'') reported outcome.
In the unbiased setting they recovered Tsirelson's CHSH bound~\cite{Tsirelson1980,CHSH1969}. For the three-outcome
CGLMP inequality~\cite{CGLMP2002} with $d=3$ they found qplex correlations with
$I_3 = 2+2/\sqrt3 \approx 3.1547$, conjecturing that this is optimal. For comparison, the local bound is 2, the
no-signalling bound is 4~\cite{PopescuRohrlich1994,Barrett2007}, and the quantum maximum is
$1+\sqrt{11/3}\approx 2.9149$~\cite{Acin2002,NPA2007,NPA2008}, proved analytically in Ref.~\cite{IoannouRosset2021}.
The standard measurements on a maximally entangled qutrit pair give the smaller value $\approx2.8729$~\cite{Acin2002}.

\paragraph{Results.}
Write $\Tq=2+2/\sqrt3$. Sections~3 and~4 prove the unbiased upper bound and identify
$\Tq$ with the maximum over $\qone$. The first-level value $3.1547$ was reported numerically
in Table~1 of Ref.~\cite{NPA2008}. Section~5 compares the attaining example with the almost-quantum set.
Section~6 establishes two properties of the full model. A biased correlation has no
level-1 moment matrix (\cref{prop:biasedQ1}), yet the score of every correlation is at most
$\Tq$ (\cref{thm:main}). Thus the two sets have the same maximum in this Bell direction,
even though the qplex model also contains points outside $\qone$.

\section{Setting}

\subsection{The CGLMP expression}
Alice and Bob each choose a setting $x,y\in\{0,1\}$ and obtain $a,b\in\mathbb{Z}_3$. Let
$Q^{xy}_k := \Pr(a-b\equiv k\mid x,y)$. We use
\begin{equation}\label{eq:I3}
I_3 = (Q^{00}_0-Q^{00}_2) + (Q^{10}_2-Q^{10}_0) + (Q^{11}_0-Q^{11}_2) + (Q^{01}_0-Q^{01}_1),
\end{equation}
which is the CGLMP expression~\cite{CGLMP2002} for $d=3$ ($I_{2233}$ in Ref.~\cite{GuptaPienaar2026}).
Ref.~\cite{GuptaPienaar2026} writes the classes as $a+b\equiv k$; with their explicit vectors this agrees with
\eqref{eq:I3} after relabelling Bob's outcomes $b\mapsto -b$. (Using their vectors with the raw labels in the sum
convention gives $Q^{xy}_k=1/3$ for all classes and $I_3=0$.) Local deterministic strategies give
$I_3\in\{2,-1,-4\}$ (30, 48 and 3 of the 81 strategies), so the local bound is 2.

\subsection{Qplexes}
Let $d=3$ and let $\Delta$ be the probability simplex in $\mathbb{R}^9$. Following Ref.~\cite{Qplex2017},
a \emph{germ} is a subset of $\Delta$ in which every pair (including a vector with itself) satisfies
$\tfrac1{12}\le p\cdot q\le\tfrac16$; a \emph{qplex} is a maximal germ. For quantum states and a SIC,
$\tr\rho\sigma = 12\,p\cdot q-1$, so the interval is $0\le\tr\rho\sigma\le1$. Bob's measurement with outcomes
$b$ is described by weights $\gamma_b\ge0$ and normalised response vectors $v_b$ with $\sum_b\gamma_b v_b=c$,
where $c=(1/9,\dots,1/9)$. The Urgleichung for $d=3$ reads
$\Pr(b\mid p) = \sum_i [4p_i-\tfrac13]\, 9\gamma_b (v_b)_i = \gamma_b(36\,p\cdot v_b-3)$.

\subsection{Bipartite qplex correlations}
In the model of Ref.~\cite{GuptaPienaar2026}, Bob holds a prior $s$ in a qplex. Alice's setting $x$ and outcome $a$ (probability $\mu^x_a$) update Bob's state
to $u^x_a$, with $\sum_a\mu^x_a u^x_a=s$ for both $x$ (no-signalling). The joint distribution is
\begin{equation}\label{eq:joint}
p(a,b\mid x,y) = \mu^x_a\,\gamma^y_b\,(36\,u^x_a\cdot v^y_b-3),
\end{equation}
where all $u^x_a$ and $v^y_b$ belong to one qplex. We call the box \emph{unbiased} if $\mu^x_a=\gamma^y_b=1/3$ and
$s=c$.

\subsection{Bloch coordinates}
Put $n(p):=\sqrt{18}\,(p-c)$, so $p\cdot q=\tfrac19+\tfrac1{18}n(p)\cdot n(q)$. The germ conditions become
\begin{equation}\label{eq:bloch}
|n|\le1,\qquad n\cdot n'\ge-\tfrac12 ,
\end{equation}
and \eqref{eq:joint} becomes $p(a,b\mid x,y)=\mu^x_a\gamma^y_b\,(1+2A^x_a\cdot B^y_b)$ with
$A^x_a=n(u^x_a)$, $B^y_b=n(v^y_b)$, $\sum_a\mu^x_aA^x_a=S:=n(s)$ and $\sum_b\gamma^y_bB^y_b=0$.
The Gram-matrix upper bounds below allow arbitrary ambient dimension and omit SIC-coordinate
non-negativity. Section~6 also relaxes Bob's norm bound. Each relaxation contains every admissible qplex realisation.

\subsection{The Gupta--Pienaar example}
With $\omega=e^{2\pi i/3}$, $h=\tfrac14[(1+\sqrt3)+i(1-\sqrt3)]$, $k=\tfrac12(-1+i)$ and
$f_w(r,s)=\tfrac19[1+\re(w_0\omega^r+w_1\omega^s)]$ on $(r,s)\in\mathbb{Z}_3^2$, the vectors
$u^0_a=f_{\omega^{-a}(1,0)}$, $u^1_a=f_{\omega^{-a}(0,1)}$, $v^0_b=f_{\omega^{-b}(h,k)}$, $v^1_b=f_{\omega^{-b}(\bar h,h)}$
form a germ. Any germ extends to a maximal germ~\cite{Qplex2017}. With unbiased weights every setting pair has $(Q_+,Q_-)=(\eta,0)$, $\eta=\tfrac12(1+1/\sqrt3)$, so
$I_3=4\eta=\Tq$.

\section{Unbiased correlations}

\begin{theorem}\label{thm:flat}
Every unbiased qplex box satisfies $I_3\le\Tq$. More generally the bound holds for any vectors $A^x_a,B^y_b$ in a
real inner-product space with $\sum_aA^x_a=\sum_bB^y_b=0$, $|A^x_a|,|B^y_b|\le1$ and
$1+2A^x_a\cdot B^y_b\ge0$, with $p(a,b|x,y)=\tfrac19(1+2A^x_a\cdot B^y_b)$.
\end{theorem}

\begin{proof}
Define complex vectors $\hat a_x=\tfrac{\sqrt2}{3}\sum_a\omega^aA^x_a$ and $\hat b_y=\tfrac{\sqrt2}{3}\sum_b\omega^bB^y_b$.
Parseval and $\sum_aA^x_a=0$ give $\|\hat a_x\|^2=\tfrac13\sum_a|A^x_a|^2\le1$, and likewise for $\hat b_y$.
A direct computation gives $\sum_k\omega^kQ^{xy}_k=\kappa_{xy}:=\langle\hat a_x,\hat b_y\rangle$ (Hermitian product,
linear in the first slot), hence
\[
Q^{xy}_k=\tfrac13\bigl(1+2\re(\omega^{-k}\kappa_{xy})\bigr),\qquad
I_3=\tfrac{2}{\sqrt3}\sum_{x,y}\re\bigl(e^{i\varphi_{xy}}\kappa_{xy}\bigr),
\]
with $\varphi_{00}=\varphi_{11}=-\pi/6$, $\varphi_{10}=5\pi/6$, $\varphi_{01}=\pi/6$. Positivity of the class sums is
$\re(\omega^m\kappa_{xy})\ge-\tfrac12$ for $m\in\mathbb{Z}_3$. Let $K\PSD0$ be the $4\times4$ Hermitian Gram matrix of
$(\hat a_0,\hat a_1,\hat b_0,\hat b_1)$, so $K_{ii}\le1$. Let $W$ be the Hermitian matrix with
$\tr(WK)=\sum_{xy}\re(e^{i\varphi_{xy}}\kappa_{xy})$, and $E_{c}$ the Hermitian matrix with
$\tr(E_cK)=\re(\omega^{m}\kappa_{xy})$ for the four constraints
$c\in\{(\hat a_0,\hat b_0,1),(\hat a_0,\hat b_1,2),(\hat a_1,\hat b_0,0),(\hat a_1,\hat b_1,1)\}$. Set $\lambda=(\sqrt3-1)/2$ and use $h,k$ from Section~2.5. In the displayed ordering of the Gram vectors,
\[
Z=\tfrac12 I_4-W-\lambda\sum_c E_c
=\tfrac12\begin{pmatrix}I_2&-D^\dagger\\-D&I_2\end{pmatrix},
\qquad
D=\begin{pmatrix}h&k\\ \bar h&h\end{pmatrix}.
\]
Indeed $D_{yx}=e^{i\varphi_{xy}}+\lambda\omega^{m_{xy}}$ for the four choices above.
The identities $|h|^2=|k|^2=1/2$ and $h^2+k\bar h=0$ give $DD^\dagger=I_2$.
It follows that $Z^2=Z$ and $\tr Z=2$. Thus $Z$ is a positive semidefinite projector of rank two.
Taking its trace pairing with $K$ gives
\[
\sum_{x,y}\re(e^{i\varphi_{xy}}\kappa_{xy})
\le\tfrac12\tr K+4\lambda/2\le 2+2\lambda=1+\sqrt3.
\]
Multiplication by $2/\sqrt3$ yields the bound. The example in Section~2.5 attains it.
\end{proof}

\begin{remark}[Scope of the Fourier bound]
The proof uses neither non-negativity of SIC coordinates nor any Alice--Alice or Bob--Bob germ condition.
Moreover, in Bloch-free form it only needs the Fourier norms $\|\hat a_x\|,\|\hat b_y\|\le1$. For biased boxes
the same certificate, applied to the augmented Fourier vectors $\tilde A_x=\sum_a\omega^a\mu^x_a(1,\sqrt2A^x_a)$ and $\tilde B_y=\sum_b\omega^b\gamma^y_b(1,\sqrt2B^y_b)$ (whose Hermitian products are the Fourier coefficients of the class distributions), yields $\tfrac{\sqrt3}{2}I_3\le1+\sqrt3+2|S|^2-\tfrac32(o_A+o_B)$, where
$o_A=\sum_x\sum_{a<a'}\mu^x_a\mu^x_{a'}(1+2A^x_a\cdot A^x_{a'})\ge0$ and $o_B$ is defined analogously.
This proves $I_3\le\Tq$ whenever $4|S|^2\le3(o_A+o_B)$, but not in general (see \cref{sec:general}).
\end{remark}

\section{Macroscopic locality}

$\qone$ is the set of no-signalling boxes admitting a level-1 NPA certificate~\cite{NPA2007,NPA2008}: a Gram
representation of $\{\mathbf{1}\}\cup\{A_{a|x}\}\cup\{B_{b|y}\}$ with
$\langle\mathbf{1},A_{a|x}\rangle=p(a|x)$, $\langle A_{a|x},A_{a'|x}\rangle=\delta_{aa'}p(a|x)$,
$\langle A_{a|x},B_{b|y}\rangle=p(ab|xy)$, and similarly for Bob. Inner products between different settings of the
same party are free. It coincides with the set of macroscopically local boxes~\cite{NavascuesWunderlich2010}.

Write $I$ and $J$ for the $3\times3$ identity and all-ones matrices.
\begin{lemma}[Gram completion for one measurement]\label{lem:pad}
If $\sum_aA_a=0$ and $A_a\cdot A_{a'}\ge-\tfrac12$ for $a\ne a'$, then
$\bigl[\tfrac29A_a\cdot A_{a'}\bigr]_{aa'}\preceq\tfrac13I-\tfrac19J$.
\end{lemma}
\begin{proof}
The difference has zero row sums and off-diagonal entries $-\tfrac19(1+2A_a\cdot A_{a'})\le0$, so it is a graph
Laplacian with non-negative weights.
\end{proof}

\begin{theorem}\label{thm:Q1member}
Every unbiased qplex box lies in $\qone$.
\end{theorem}
\begin{proof}
Use a common unit vector $e_0$ orthogonal to the Bloch space. Take $v_{\mathbf{1}}=e_0$ and
$v_{A_{a|x}}=e_0/3+(\sqrt2/3)A^x_a+\pi^x_a$,
$v_{B_{b|y}}=e_0/3+(\sqrt2/3)B^y_b+\pi'^y_b$, where the padding vectors $\pi$ live in mutually
orthogonal fresh subspaces (one per measurement), sum to zero within each measurement, and have Gram matrix
$\tfrac13I-\tfrac19J-[\tfrac29A_a\cdot A_{a'}]$ (positive semidefinite by \cref{lem:pad}). Then
$\langle v_A,v_B\rangle=\tfrac19+\tfrac29A\cdot B=p(ab|xy)$ and all level-1 constraints hold. For the Gupta--Pienaar example the padding vanishes because
distinct outcomes of each measurement have Bloch inner product $-1/2$ and each vector has unit norm.
\end{proof}

\begin{theorem}\label{thm:Q1max}
$\max_{\qone}I_3=\Tq$.
\end{theorem}
\begin{proof}
The example in Section~2.5 and \cref{thm:Q1member} give the lower bound. For the upper bound, $I_3$ is invariant under the joint relabelling
$(a,b)\mapsto(a+1,b+1)$, which preserves $\qone$; averaging over this $\mathbb{Z}_3$ action makes all marginals
uniform without changing $I_3$. For uniform marginals, a level-1 certificate yields centred vectors
$\alpha_{a|x}=v_{A_{a|x}}-\tfrac13v_{\mathbf{1}}$ with $\sum_a\alpha_{a|x}=0$, within-setting Gram
$\tfrac13I-\tfrac19J$ and $\alpha\cdot\beta=p-\tfrac19$. Here a real Gram representation suffices by taking real parts. Setting $A:=\tfrac3{\sqrt2}\alpha$ gives $|A|=1$ and
$p=\tfrac19(1+2A\cdot B)\ge0$, so \cref{thm:flat} applies.
\end{proof}

\begin{remark}
Table~1 of Ref.~\cite{NPA2008} gives the first-level value $3.1547$. If one omits
$p(ab|xy)\ge0$ the level-1 SDP returns 4 (cf.\ the footnote in Ref.~\cite{NavascuesWunderlich2010}).
Our SDP reproduces $3.1547005$ with positivity imposed.
\end{remark}

\begin{corollary}
The optimal unbiased qplex value of $I_3$ equals the macroscopic-locality bound. Macroscopic locality therefore
cannot exclude these qplex correlations.
\end{corollary}

\section{Almost-quantum correlations}

\begin{proposition}\label{prop:aq}
The Gupta--Pienaar box is not in the almost-quantum set $\qab$~\cite{AlmostQuantum2015}.
\end{proposition}
\begin{proof}
The level-$1+AB$ sum-of-squares certificate of Ioannou and Rosset~\cite{IoannouRosset2021}
bounds $I_3$ by $1+\sqrt{11/3}$ on $\qab$, below the value of the example.
This agrees with the level-$1+AB$ result of Ref.~\cite{NPA2008}. The accompanying
discovery code also rounds a numerical dual witness to rationals and checks positivity
by exact elimination~\cite{PeyrlParrilo2008}, giving the sufficient bound $I_3<2.915$.
\end{proof}

\begin{remark}[Numerical visibility]
Let $p_v=v p_{\rm GP}+(1-v)\mathbf{u}$, where $\mathbf{u}(a,b|x,y)=1/9$.
The numerical visibility threshold for membership in $\qab$ is approximately $0.888889$,
close to $8/9$. The CGLMP upper bound alone gives $v\le(1+\sqrt{11/3})/\Tq\approx0.924$.
The exact visibility remains open here.
\end{remark}

\section{General weights}\label{sec:general}

\begin{proposition}\label{prop:biasedQ1}
There is a biased, exactly no-signalling qplex correlation outside $\qone$.
One rational example has $I_3=391/144$.
\end{proposition}
\begin{proof}
Take $\mu^0=\mu^1=(3/5,1/5,1/5)$, $\gamma^0=(1/3,1/3,1/3)$ and
$\gamma^1=(3/8,5/16,5/16)$. Appendix~\ref{app:biased} specifies the twelve rational
vectors. Their rows sum to one, their coordinates are positive, and their pairwise inner
products, including self-pairs, satisfy
\[
 \frac1{12}<\frac{13547}{162000}\le u\cdot v\le\frac{2681}{16200}<\frac16
 \qquad (u,v\text{ any two of the twelve vectors}).
\]
Their barycentres satisfy both equalities of Section~2, so \eqref{eq:joint} defines an
exactly no-signalling qplex box.

A level-1 representation would yield a positive semidefinite Gram matrix $\Gamma$
for the unit vector and the first two outcomes of each of the four measurements.
Only the blocks between distinct settings of one party are unspecified.
The integer matrix $H$ in Appendix~\ref{app:biased} has zeros there and is positive definite.
Substitution of the fixed entries gives $\tr(H\Gamma)=-1189/900$, contradicting
$\tr(H\Gamma)\ge0$. Direct evaluation of \eqref{eq:I3} gives the stated score.
\end{proof}
The example was obtained by rationalising and perturbing a numerical candidate from the
first draft. The earlier candidate had numerical white-noise visibility near $0.93188$. The exact separator applies to the rational coordinates in the appendix.
Since $\qone$ membership fails in the full model, a separate upper-bound argument is required.

\begin{theorem}[General CGLMP bound]\label{thm:main}
For all weights $\mu^x_a,\gamma^y_b$, all priors $s$ and all germs, the qplex correlations \eqref{eq:joint} satisfy
$I_3\le\Tq=2+2/\sqrt3$. The bound is attained. Equality requires $\mu^0=\mu^1=(1/3,1/3,1/3)$.
\end{theorem}

The proof has four steps.

\paragraph{Step 1. Reduction to Alice's weights.}
Introduce $e_0=(1,0)$ and homogenised vectors $\tilde a_{x,a}=\mu^x_a(1,\sqrt2A^x_a)$,
$\tilde b_{y,b}=\gamma^y_b(1,\sqrt2B^y_b)$. Then $p(a,b|x,y)=\tilde a_{x,a}\cdot\tilde b_{y,b}$.
All twelve vectors $u^x_a,v^y_b$ lie in one qplex, so $n\cdot n'\ge-\tfrac12$ holds for every pair of them,
including pairs from different settings or parties. Hence
$w\cdot w'\ge0$ for every pair of homogenised outcome vectors, Bob's normalisation becomes $\sum_b\tilde b_{y,b}=e_0$, no-signalling becomes
$\sum_a\tilde a_{0,a}=\sum_a\tilde a_{1,a}$, and Alice's purity is $|\tilde a_{x,a}|^2\le3(\mu^x_a)^2$.
Dropping Bob's norm bound gives an upper bound for every fixed $\mu=(\mu^0,\mu^1)$.
With $X$ the $13\times13$ Gram matrix of $(e_0,\tilde a,\tilde b)$, define
\begin{equation}\label{eq:Psi}
\begin{aligned}
\Psi(\mu)=\max\ &\textstyle\sum_{xy}\sum_a\bigl(X_{\tilde a_{x,a},\tilde b_{y,a-g_{xy}}}-X_{\tilde a_{x,a},\tilde b_{y,a-l_{xy}}}\bigr)\\
\text{s.t. }&X\PSD0,\ X_{00}=1,\ X_{\tilde a_i,0}=\mu_i,\ X_{\tilde a_i\tilde a_i}\le3\mu_i^2,\ X_{ij}\ge0\ (1\le i<j),\\
&X(\chi_{B_y}-e_0)=0\ (y=0,1),\quad X(\chi_{A_0}-\chi_{A_1})=0,
\end{aligned}
\end{equation}
where $(g_{xy},l_{xy})$ are the plus and minus classes in \eqref{eq:I3}. In the last line,
$\chi_{A_x},\chi_{B_y}\in\{0,1\}^{13}$ indicate the three rows of $X$ belonging to
$\tilde a_{x,\cdot}$ and $\tilde b_{y,\cdot}$, and $e_0$ also denotes the indicator of the first row. Every qplex configuration with Alice
weights $\mu$ is feasible, so $I_3\le\Psi(\mu)$. The reduction also covers zero weights
because it only multiplies by the weights and never divides by them.

\paragraph{Step 2. A parameter-free dual.}
In \eqref{eq:Psi}, $\mu$ appears only on right-hand sides. Hence the dual feasible set $\mathcal D$ (vectors of
multipliers $y$ with $Z(y):=\mathcal A(y)-C\PSD0$ and sign constraints) does not depend on $\mu$, and weak duality gives,
for every $y\in\mathcal D$ and every $\mu$,
\begin{equation}\label{eq:fy}
\Psi(\mu)\le f_y(\mu):=y_{00}+\sum_i y_{A,i}\,\mu_i+3\sum_i z_{A,i}\,\mu_i^2,\qquad z_{A,i}\ge0 .
\end{equation}
For completeness, let $E_{ij}=(e_i e_j^\top+e_j e_i^\top)/2$ in $\mathbb R^{13}$,
let $C$ be the symmetric objective matrix from \eqref{eq:Psi}, and let $t_1,t_2,t_3$ be the vectors
for which the last line of \eqref{eq:Psi} requires $X t_j=0$. The slack used by the verifier is
\[
Z(y)=y_{00}E_{00}+\sum_{i=1}^6(y_{A,i}E_{i0}+z_{A,i}E_{ii})
-\sum_{1\le i<j\le12}u_{ij}E_{ij}
-\sum_{j=1}^3 (w_jt_j^\top+t_jw_j^\top)/2-C .
\]
Here $u_{ij},z_{A,i}\ge0$ and $w_j\in\mathbb R^{13}$ is free.
Pairing with a feasible $X$ eliminates the three equality terms and gives \eqref{eq:fy}.
Each $f_y$ is convex in $\mu$, so its maximum on any parameter box is attained at a vertex.

\begin{remark}[Behaviour near the uniform point]
The dual coefficients can be chosen separately for each weight vector. At the uniform
point the constructed slack is singular. We use a quadratic family of multipliers to
control the directions in its kernel as the weights vary.
\end{remark}

\paragraph{Step 3. Outer region.}
Write $p=(\mu^0_0,\mu^0_1,\mu^1_0,\mu^1_1)$ for the four independent weights.
Their domain is the subset of the unit four-cube satisfying
$p_0+p_1\le1$ and $p_2+p_3\le1$. For each of \NBOX{} closed dyadic boxes, the supplement supplies a
rational $y$. The verifier checks $u,z_A\ge0$, positive pivots for $Z(y)$, and
$f_y(\mu_v)\le3154700538/10^9<\Tq$ at all 16 vertices. The largest vertex value is about
$3.1546986$, roughly $2\cdot10^{-6}$ below $\Tq$. Convexity extends the bound to the whole box.

The verifier reconstructs the subdivision from a $4^4$ grid, splitting the first longest edge
at its midpoint. It checks that the boxes cover the interior of the weight domain apart from
16 further boxes, each contained in $\|\mu-\bar\mu\|\le1/125$. Here $\bar\mu=(1/3,\ldots,1/3)$
and the norm uses all six weights. Containment follows from the exact squared distances of the
vertices. The finite union of these closed boxes also contains the boundary by closure.

\paragraph{Step 4. A local certificate at the optimum.}
Let $u=(1,-1,0)$, $v=(1,1,-2)$ and use the four tangent vectors
$(u,0),(v,0),(0,u),(0,v)$ in that order. Write $\mu(\delta)=\bar\mu+\sum_i\delta_i\tau_i$,
where $\tau_i$ denotes the $i$th vector. Take the quadratic family
\[
y(\delta)=y_0+Y_1\delta+\textstyle\sum_{i,j}\delta_i\delta_jY_{2,ij}.
\]
Because $\mathcal D$ is $\mu$-independent and $Z$ is affine in $y$, $Z(y(\delta))=\mathcal Z_0+\mathcal Z_1(\delta)+\mathcal Z_2(\delta)$ is
exactly quadratic in $\delta$. The certificate supplies an invertible matrix $G$ over $\mathbb Q(\sqrt3)$ partitioned into eight and five columns. Write $\hat{\mathcal Z}=G^\top\mathcal ZG$ in $(pp,pk,kk)$ blocks.

\begin{lemma}[Scaled congruence]\label{lem:scaled}
Suppose $\hat{\mathcal Z}_0=\operatorname{diag}(Z_{0p},0)$ and $\hat{\mathcal Z}_{1,kk}\equiv0$.
Form the block matrices $\widetilde W=[\hat{\mathcal Z}_{1,i,pk}]_i$ and
$\widetilde K=[\hat{\mathcal Z}_{2,ij,kk}]_{ij}$. Assume
\[
 \begin{pmatrix} Z_{0p}&\widetilde W\\ \widetilde W^\top&\widetilde K\end{pmatrix}
 \succeq m_0I.
\] Let
$e_1,e_2,b_2$ bound $\|\hat{\mathcal Z}_{1,pp}(d)\|$, $\|\hat{\mathcal Z}_{2,pp}(d)\|$ and
$\|\hat{\mathcal Z}_{2,pk}(d)\|$, respectively, for $|d|=1$. Then
$Z(y(\delta))\succ0$ whenever $0<|\delta|=\rho$ and $m_0>\rho(e_1+b_2)+\rho^2e_2$.
\end{lemma}
\begin{proof}
With $S_\rho=\operatorname{diag}(I_8,\rho^{-1}I_5)$ and $\delta=\rho d$,
$S_\rho^\top\hat{\mathcal Z}(\rho d)S_\rho=M_0(d)+\rho M_1(d)+\rho^2M_2(d)$ with
$M_0(d)=\begin{psmallmatrix}Z_{0p}&\hat{\mathcal Z}_{1,pk}(d)\\ \cdot&\hat{\mathcal Z}_{2,kk}(d)\end{psmallmatrix}$.
Testing the hypothesis matrix on vectors $(x,d\otimes z)$ shows $M_0(d)\succeq m_0I$ for $|d|=1$, while
$\|M_1(d)\|\le e_1+b_2$ and $\|M_2(d)\|\le e_2$.
\end{proof}

Let $R=1/125$. The file \texttt{local.json} supplies $G,y_0,Y_1,Y_2$.
The verifier checks that $G$ is invertible (via $G^\top G\succeq I/2$), the kernel blocks in
\cref{lem:scaled} vanish, and the unshifted $28\times28$ block matrix has margin $m_0=7/50$.
For $f(\delta)=f_{y(\delta)}(\mu(\delta))$, it checks $f(0)=\Tq$, zero linear term,
and $-Q\succeq cI$ with $c=7/25$, where $\delta^\top Q\delta$ is the quadratic term.
The cubic and quartic terms are bounded in magnitude by $r_3|\delta|^3$ and $r_4|\delta|^4$.
The following upper estimates all hold in exact arithmetic:
\[
 e_1\le11,\qquad e_2\le5,\qquad b_2\le1/50,\qquad r_3\le26,\qquad r_4\le16.
\]
In particular, $R(11+1/50)+5R^2<7/50$ and $26R+16R^2<7/25$. The second inequality gives $f(\delta)<\Tq$ whenever $0<|\delta|\le R$.
Thus $Z(y(\delta))\succeq0$ and $f(\delta)\le\Tq$ for $|\delta|\le R$.
For each sign-constrained multiplier that vanishes at the centre, the linear term vanishes and
the quadratic coefficient matrix is positive definite. For each remaining multiplier, its
positive constant exceeds the sum of the worst-case linear and quadratic variations on this ball.
These tests establish dual feasibility throughout the ball.

For the two eigenvalue margins the verifier rounds each entry of the exact symmetric matrix
to a rational grid and confirms an absolute error at most $10^{-7}$ by exact comparisons in
$\mathbb Q(\sqrt3)$. For an $n\times n$ matrix this bounds the operator norm of the error by
$n\,10^{-7}$. Rational elimination on the shifted matrix proves the required lower bound.
For the remainder estimates, entrywise rational enclosures give Frobenius bounds on the
coefficient tensors. Finally,
\[
 \|\mu-\bar\mu\|^2=2\delta_1^2+6\delta_2^2+2\delta_3^2+6\delta_4^2\ge 2|\delta|^2 .
\]
The local certificate therefore covers the entire weight-space ball of radius $\sqrt2/125$,
including the 16 inner boxes of Step~3.

\begin{proof}[Proof of \cref{thm:main}]
Step~1 gives $I_3\le\Psi(\mu)$. The outer boxes have a strict upper bound below $\Tq$.
In the inner boxes, Step~4 gives a strict upper bound whenever $\delta\ne0$.
Thus equality requires uniform Alice weights. The example in Section~2.5 attains equality.
\end{proof}

\paragraph{Numerical exploration.}
The original search also found nearly two-outcome weight vectors with scores around $3.10610$,
for instance with $\mu^0\approx(0.4445,0.5555,0)$. These observations were used to explore
the weight domain before constructing the global certificate.

\section{Discussion}

For the Gupta--Pienaar single-system construction, qplex geometry and macroscopic locality
give the same maximum of $I_3$. The rational example in \cref{prop:biasedQ1} clarifies the
scope of that equality. Its CGLMP value is $391/144\approx2.7153$, below both maxima,
but its fixed marginal and joint probabilities preclude a level-1 Gram matrix. The CGLMP
score alone cannot detect this failure of macroscopic locality.

The Gupta--Pienaar maximiser also distinguishes $\qone$ from $\qab$ in this scenario.
Whether local orthogonality~\cite{Fritz2013} or information causality~\cite{Pawlowski2009}
excludes this particular box is a separate question. Multi-outcome information-causality
criteria~\cite{Gachechiladze2022} provide possible tests.
Table~1 of Ref.~\cite{NPA2008} gives a first-level bound of approximately $3.2126$ for
$d=4$. Whether the corresponding qplex model attains it remains open here. So does the
construction of a composite-system qplex realising the correlations considered in this paper.

\section*{Contributions}
Yusuf Kaan Öğüt posed the optimality problem, directed the investigation, chose to include
arbitrary weights, and takes responsibility for the manuscript. Claude Opus 5.5 (Anthropic)
produced the first draft, the initial proofs, the computational search and its code under his direction.
OpenAI Codex (GPT-6.2 Astra) reviewed the arguments, implemented and ran a separate checker
for the fixed certificates using the original quadratic-field arithmetic, derived the projector
form of the unbiased certificate, constructed the rational biased example and its separating
matrix, and revised the manuscript. Both named systems are credited as AI contributors.

\section*{Code and data availability}
The supplement includes the fixed witnesses, a verifier using the Python standard library,
and the original search programs. Numerical discovery used NumPy~\cite{NumPy2020},
SciPy~\cite{SciPy2020} and CVXPY~\cite{CVXPY2016} with Clarabel~\cite{Clarabel2024}
and SCS~\cite{SCS2016}.

\appendix
\section{Checking the certificates}\label{app:repro}
From the root of the supplement, run
\begin{center}\texttt{python3 certificate/verify.py}\end{center}
with Python 3.9 or later. The verifier evaluates the example attaining $\Tq$, checks the rational
separator in Appendix~\ref{app:biased}, reconstructs each outer slack and all 16 vertex
values, and verifies the local polynomial certificate. The checker rebuilds the matrices
from the Bell coefficients and the equalities in \eqref{eq:Psi}. The data files in
\texttt{certificate/} are \texttt{biased.json} (twelve states and a separating matrix),
\texttt{boxes.json} (dyadic leaves and rational dual multipliers) and \texttt{local.json}
(change of basis and quadratic dual coefficients). They contain rational numbers or pairs of
rationals representing $a+b\sqrt3$, so no optimisation solver is needed.


\clearpage
\section{A rational qplex box outside macroscopic locality}\label{app:biased}
Each row below gives a common denominator $d_r$ and nine numerators. Coordinates
are in lexicographic order on $\mathbb Z_3^2$. The weights are those of \cref{prop:biasedQ1}.
\begin{center}
\small
\begin{tabular}{@{}lr@{\hspace{1.5em}}rrrrrrrrr@{}}
\toprule
State & $d_r$ & \multicolumn{9}{c}{Coordinate numerators}\\
\midrule
$u^{0}_{0}$ & 180 & 8 & 10 & 21 & 19 & 11 & 7 & 40 & 49 & 15 \\
$u^{0}_{1}$ & 180 & 46 & 39 & 7 & 13 & 6 & 24 & 5 & 16 & 24 \\
$u^{0}_{2}$ & 180 & 3 & 25 & 7 & 46 & 37 & 26 & 15 & 7 & 14 \\
\addlinespace[2.5pt]
$u^{1}_{0}$ & 180 & 12 & 26 & 2 & 33 & 9 & 12 & 29 & 47 & 10 \\
$u^{1}_{1}$ & 180 & 26 & 5 & 18 & 8 & 19 & 10 & 42 & 8 & 44 \\
$u^{1}_{2}$ & 180 & 11 & 11 & 53 & 9 & 30 & 25 & 11 & 21 & 9 \\
\addlinespace[2.5pt]
$v^{0}_{0}$ & 180 & 16 & 6 & 53 & 2 & 23 & 22 & 17 & 26 & 15 \\
$v^{0}_{1}$ & 180 & 39 & 43 & 5 & 22 & 8 & 26 & 1 & 22 & 14 \\
$v^{0}_{2}$ & 180 & 5 & 11 & 2 & 36 & 29 & 12 & 42 & 12 & 31 \\
\addlinespace[2.5pt]
$v^{1}_{0}$ & 180 & 9 & 17 & 9 & 28 & 12 & 9 & 37 & 45 & 14 \\
$v^{1}_{1}$ & 180 & 46 & 25 & 16 & 3 & 10 & 20 & 18 & 4 & 38 \\
$v^{1}_{2}$ & 900 & 36 & 93 & 186 & 137 & 198 & 166 & 8 & 30 & 46 \\
\addlinespace[2.5pt]
\bottomrule
\end{tabular}
\end{center}
Index the moment matrix by
\[
 (1,A_{0|0},A_{1|0},A_{0|1},A_{1|1},B_{0|0},B_{1|0},B_{0|1},B_{1|1}).
\]
The cell involving $1$ and an outcome, as well as that outcome's diagonal entry, equals
its marginal. Different outcomes of the same setting have zero inner product.
Each Alice--Bob entry is the corresponding joint probability from \eqref{eq:joint}.
The only unknown entries connect different settings of one party. At these entries the
symmetric matrix
\[
H=\begin{pmatrix}
32 & 55 & -31 & -104 & -27 & -68 & 73 & 16 & -1\\
55 & 296 & -163 & 0 & 0 & -153 & 165 & -298 & 19\\
-31 & -163 & 94 & 0 & 0 & 85 & -92 & 166 & -11\\
-104 & 0 & 0 & 601 & 155 & 223 & -241 & -417 & 27\\
-27 & 0 & 0 & 155 & 43 & 58 & -63 & -108 & 7\\
-68 & -153 & 85 & 223 & 58 & 166 & -176 & 0 & 0\\
73 & 165 & -92 & -241 & -63 & -176 & 194 & 0 & 0\\
16 & -298 & 166 & -417 & -108 & 0 & 0 & 595 & -39\\
-1 & 19 & -11 & 27 & 7 & 0 & 0 & -39 & 6
\end{pmatrix}
\]
vanishes. Its successive leading principal determinants are
\[
\begin{gathered}
32,\quad 6447,\quad 27184,\quad 2763504,\quad 8335577,\\
31460205,\quad 144732823,\quad 660682302,\quad 2199990491 .
\end{gathered}
\]
Sylvester's criterion gives $H\succ0$. Using the stated marginals and joint probabilities
in the trace pairing gives $-1189/900$ for every Hermitian completion. Since two positive
semidefinite matrices have non-negative trace pairing, a level-1 completion cannot exist.


\clearpage
\bibliographystyle{unsrt}
\bibliography{refs}

\end{document}
